\documentclass[10pt,a4paper]{amsart}
\usepackage{amsmath,amssymb,amsthm,graphicx}
\usepackage{verbatim}
\newenvironment{reference}{\comment}{\endcomment}
%\newenvironment{reference}{}{}
\newenvironment{slogan}{\comment}{\endcomment}
\newenvironment{history}{\comment}{\endcomment}

% For commutative diagrams we use Xy-pic
\usepackage[all]{xy}

% We use 2cell for 2-commutative diagrams.
\xyoption{2cell}
% Theorem environments.
%
\theoremstyle{plain}
\newtheorem{theorem}[subsection]{Theorem}
\newtheorem{proposition}[subsection]{Proposition}
\newtheorem{lemma}[subsection]{Lemma}

\theoremstyle{definition}
\newtheorem{definition}[subsection]{Definition}
\newtheorem{example}[subsection]{Example}
\newtheorem{exercise}[subsection]{Exercise}
\newtheorem{situation}[subsection]{Situation}

\theoremstyle{remark}
\newtheorem{remark}[subsection]{Remark}
\newtheorem{remarks}[subsection]{Remarks}

\numberwithin{equation}{subsection}

% Macros
%
\def\lim{\mathop{\mathrm{lim}}\nolimits}
\def\colim{\mathop{\mathrm{colim}}\nolimits}
\def\Spec{\mathop{\mathrm{Spec}}}
\def\Hom{\mathop{\mathrm{Hom}}\nolimits}
\def\Ext{\mathop{\mathrm{Ext}}\nolimits}
\def\SheafHom{\mathop{\mathcal{H}\!\mathit{om}}\nolimits}
\def\SheafExt{\mathop{\mathcal{E}\!\mathit{xt}}\nolimits}
\def\Sch{\mathit{Sch}}
\def\Mor{\mathop{\mathrm{Mor}}\nolimits}
\def\Ob{\mathop{\mathrm{Ob}}\nolimits}
\def\Sh{\mathop{\mathit{Sh}}\nolimits}
\def\NL{\mathop{N\!L}\nolimits}
\def\CH{\mathop{\mathrm{CH}}\nolimits}
\def\proetale{{pro\text{-}\acute{e}tale}}
\def\etale{{\acute{e}tale}}
\def\QCoh{\mathit{QCoh}}
\def\Ker{\mathop{\mathrm{Ker}}}
\def\Im{\mathop{\mathrm{Im}}}
\def\Coker{\mathop{\mathrm{Coker}}}
\def\Coim{\mathop{\mathrm{Coim}}}

% Boxtimes
%
\DeclareMathSymbol{\boxtimes}{\mathbin}{AMSa}{"02}

%
% Macros for moduli stacks/spaces
%
\def\QCohstack{\mathcal{QC}\!\mathit{oh}}
\def\Cohstack{\mathcal{C}\!\mathit{oh}}
\def\Spacesstack{\mathcal{S}\!\mathit{paces}}
\def\Quotfunctor{\mathrm{Quot}}
\def\Hilbfunctor{\mathrm{Hilb}}
\def\Curvesstack{\mathcal{C}\!\mathit{urves}}
\def\Polarizedstack{\mathcal{P}\!\mathit{olarized}}
\def\Complexesstack{\mathcal{C}\!\mathit{omplexes}}
% \Pic is the operator that assigns to X its picard group, usage \Pic(X)
% \Picardstack_{X/B} denotes the Picard stack of X over B
% \Picardfunctor_{X/B} denotes the Picard functor of X over B
\def\Pic{\mathop{\mathrm{Pic}}\nolimits}
\def\Picardstack{\mathcal{P}\!\mathit{ic}}
\def\Picardfunctor{\mathrm{Pic}}
\def\Deformationcategory{\mathcal{D}\!\mathit{ef}}
\usepackage{cleveref}
\crefname{originaldefinition}{original definition}{original definitions}
\Crefname{originaldefinition}{Original definition}{Original definitions}
\makeatletter
\newlabel{categories-definition-initial}{{09WP}{source}{}{sourcepointer-1}{}}
\newlabel{categories-definition-initial@cref}{{[originaldefinition][0][]09WP}{[1][0][]source}{}{}{}}
\makeatother
\crefname{originallemma}{original lemma}{original lemmas}
\Crefname{originallemma}{Original lemma}{Original lemmas}
\makeatletter
\newlabel{categories-lemma-cofinal}{{04E7}{source}{}{sourcepointer-2}{}}
\newlabel{categories-lemma-cofinal@cref}{{[originallemma][0][]04E7}{[1][0][]source}{}{}{}}
\makeatother
\crefname{originalreference}{original reference}{original references}
\Crefname{originalreference}{Original reference}{Original references}
\makeatletter
\newlabel{lemma-presheaf-value-weakly-contractible}{{0F63}{source}{}{sourcepointer-3}{}}
\newlabel{lemma-presheaf-value-weakly-contractible@cref}{{[originalreference][0][]0F63}{[1][0][]source}{}{}{}}
\makeatother
\crefname{originalsection}{original section}{original sections}
\Crefname{originalsection}{Original section}{Original sections}
\makeatletter
\newlabel{sites-section-continuous-functors}{{00WU}{source}{}{sourcepointer-4}{}}
\newlabel{sites-section-continuous-functors@cref}{{[originalsection][0][]00WU}{[1][0][]source}{}{}{}}
\makeatother
\crefname{originalsection}{original section}{original sections}
\Crefname{originalsection}{Original section}{Original sections}
\makeatletter
\newlabel{sites-section-morphism-sites}{{00X0}{source}{}{sourcepointer-5}{}}
\newlabel{sites-section-morphism-sites@cref}{{[originalsection][0][]00X0}{[1][0][]source}{}{}{}}
\makeatother
\title{The pullback of a sheaf of sets on the pro-etale site evaluated on an affine weakly contractible object}
\author{The Stacks Project Authors}
\date{stacks-project:proetale.tex:local-snapshot-20261012; GNU FDL 1.2 or later}
\begin{document}
\maketitle
\noindent\textit{Editorial note.} All delivered mathematical text is the author's own, copied from the pinned Stacks Project chapter file proetale.tex; no statement or proof step was rewritten, reordered or removed.\par\medskip
\noindent\textit{Editorial note.} The counted proof is the author's complete argument: the description of the pullback as the sheafification of the presheaf of colimits, the initiality argument showing that this presheaf turns finite disjoint unions into products, and the closing application of the agreement lemma for affine weakly contractible objects.\par\medskip
\noindent\textit{Editorial note.} The label lemma-small-pullback-weakly-contractible names the counted statement, and the proof environment immediately following it in the source is delivered verbatim as its proof.\par\medskip
\noindent\textit{Editorial note.} The counted statement is carried by the author's own lemma environment in the chapter source.\par\medskip
\noindent\textit{Editorial note.} Pre-open gate: python operations/agent-authoring/tools/precard\_gate.py --source proetale.tex --label lemma-small-pullback-weakly-contractible --start 3499 --end 3538 --local-id 3592 returned RESULT=OK (L3592 free, occupancy FREE, public interval clear). Candidate taken from sources/stacks-proetale/candidates-halfA.json.\par\medskip
\noindent\textit{Editorial note.} The statement and proof are copied from the pinned chapter file; the cited inputs are kept as external pointers with their canonical Stacks tags.\par\medskip
\noindent\textit{Allowed standard prerequisites (editorial).}\begin{itemize}
\item Description of the pullback as the sheafification of a colimit presheaf. Sites, the sections describing the pullback along a morphism of sites as the sheafification of the presheaf of colimits over the morphisms to objects of the source site. Supplies the description of the pullback functor as the sheafification of the presheaf u\_p F of colimits.
\item Agreement of the sheafification with the presheaf on affine weakly contractible objects. Same chapter, the lemma that a presheaf of sets sending finite disjoint unions to products agrees with its sheafification on affine weakly contractible objects. Supplies the reduction of the claim to showing that the presheaf u\_p F sends finite disjoint unions to products.
\item Initiality and cofinality for a product of categories over a disjoint union. Categories, the definition of an initial functor and the lemma that the induced functor on opposite categories is cofinal, hence computes the same colimit. Supplies the identification of the colimit over morphisms from a disjoint union with the colimit over the product of the categories of morphisms from the pieces.
\end{itemize}\medskip
\ifdefined\hypertarget\hypertarget{sourcepointer-1}{}\fi
\noindent\textit{Source reference.} Original definition 09WP: sources/stacks-chow/tags.txt, source lines 12502--12502. Prerequisite pointer invoked at proof line(s) 3531; it names a standard fact of a sibling chapter that the counted argument uses, and every step of the counted proof built on it is retained verbatim in the delivered body.\par\medskip
\ifdefined\hypertarget\hypertarget{sourcepointer-2}{}\fi
\noindent\textit{Source reference.} Original lemma 04E7: sources/stacks-chow/tags.txt, source lines 5753--5753. Prerequisite pointer invoked at proof line(s) 3533; it names a standard fact of a sibling chapter that the counted argument uses, and every step of the counted proof built on it is retained verbatim in the delivered body.\par\medskip
\ifdefined\hypertarget\hypertarget{sourcepointer-3}{}\fi
\noindent\textit{Source reference.} Original reference 0F63: sources/stacks-proetale/pinned-original-source.tex, source lines 3465--3465. Prerequisite pointer invoked at proof line(s) 3521; it names a standard fact of the same chapter that the counted argument uses, and every step of the counted proof built on it is retained verbatim in the delivered body.\par\medskip
\ifdefined\hypertarget\hypertarget{sourcepointer-4}{}\fi
\noindent\textit{Source reference.} Original section 00WU: sources/stacks-chow/tags.txt, source lines 1196--1196. Prerequisite pointer invoked at proof line(s) 3518; it names a standard fact of a sibling chapter that the counted argument uses, and every step of the counted proof built on it is retained verbatim in the delivered body.\par\medskip
\ifdefined\hypertarget\hypertarget{sourcepointer-5}{}\fi
\noindent\textit{Source reference.} Original section 00X0: sources/stacks-chow/tags.txt, source lines 1202--1202. Prerequisite pointer invoked at proof line(s) 3517; it names a standard fact of a sibling chapter that the counted argument uses, and every step of the counted proof built on it is retained verbatim in the delivered body.\par\medskip
\begin{lemma}
\label{lemma-small-pullback-weakly-contractible}
Let $f : X \to Y$ be a morphism of schemes. Let $\mathcal{F}$ be a sheaf
of sets on $Y_\proetale$. If $W$ is an affine weakly contractible
object of $X_\proetale$, then
$$
f_{small}^{-1}\mathcal{F}(W) = \colim_{W \to V} \mathcal{F}(V)
$$
where the colimit is over morphisms $W \to V$ over $Y$
with $V \in Y_\proetale$.
\end{lemma}

\begin{proof}
Recall that $f_{small}^{-1}\mathcal{F}$ is the sheaf associated to the
presheaf
$$
u_p\mathcal{F} : U \mapsto \colim_{U \to V} \mathcal{F}(V)
$$
on $X_\etale$, see Sites, Sections \ref{sites-section-morphism-sites} and
\ref{sites-section-continuous-functors}; we've suppressed from the notation
that the colimit is over the opposite of the category
$\{U \to V, V \in Y_\proetale\}$. By
Lemma \ref{lemma-presheaf-value-weakly-contractible}
it suffices to prove that $u_p\mathcal{F}$ sends finite disjoint unions
to products.
Suppose that $U = U_1 \amalg U_2$ is a disjoint union of open and closed
subschemes. There is a functor
$$
\{U_1 \to V_1\} \times \{U_2 \to V_2\} \longrightarrow
\{U \to V\},\quad
(U_1 \to V_1, U_2 \to V_2) \longmapsto (U \to V_1 \amalg V_2)
$$
which is initial (Categories, Definition \ref{categories-definition-initial}).
Hence the corresponding functor on opposite categories is cofinal and
by Categories, Lemma \ref{categories-lemma-cofinal} we see
that $u_p\mathcal{F}$ on $U$ is the colimit of the values
$\mathcal{F}(V_1 \amalg V_2)$ over the product category.
Since $\mathcal{F}$ is a sheaf it sends disjoint unions to products and
we conclude $u_p\mathcal{F}$ does too.
\end{proof}

\end{document}
